% Generated draft: requires mathematical and source review.
\documentclass[11pt]{article}
\usepackage[margin=25mm]{geometry}
\usepackage{fontspec}
\setmainfont{Latin Modern Roman}
\usepackage{amsmath,amssymb,mathtools,unicode-math}
\setmathfont{Latin Modern Math}
\usepackage{enumitem,needspace,xurl,hyperref}
\setlength{\parindent}{0pt}
\setlength{\parskip}{5pt}
\setlength{\emergencystretch}{3em}
\newcommand{\N}{\mathbb{N}}
\newcommand{\Z}{\mathbb{Z}}
\newcommand{\Q}{\mathbb{Q}}
\newcommand{\R}{\mathbb{R}}
\newcommand{\C}{\mathbb{C}}
\newcommand{\F}{\mathbb{F}}
\newcommand{\A}{\mathbb{A}}
\newcommand{\PP}{\mathbb{P}}
\newcommand{\E}{\mathbb{E}}
\newcommand{\eps}{\varepsilon}
\newcommand{\dd}{\,\mathrm{d}}
\begin{document}
\section*{UnsolvedMath ID: 701735}
\textbf{Source problem code:} catalogue entry
\subsection{Definitions and mathematical statement}
A group \(G\) is \emph{boundedly generated} if there exist an integer \(k\geq 1\) and elements \(g_1,\ldots,g_k\in G\) such that
\[
G=\langle g_1\rangle\cdots\langle g_k\rangle
=\{g_1^{a_1}\cdots g_k^{a_k}:a_1,\ldots,a_k\in\Z\}.
\]
Here the product is a setwise product of cyclic subgroups; the cyclic subgroups may be finite or infinite. A group \(G\) is \emph{residually finite} if for every \(g\in G\setminus\{1\}\) there are a finite group \(Q\) and a homomorphism \(\varphi:G\to Q\) with \(\varphi(g)\neq 1\). It is \emph{linear over a field of characteristic zero} if there exist a field \(K\) with \(\operatorname{char}(K)=0\), an integer \(n\geq 1\), and an injective homomorphism \(G\hookrightarrow\operatorname{GL}_n(K)\).

\textbf{Question.} Does every boundedly generated residually finite group admit such a faithful finite-dimensional representation over some field of characteristic zero? Both the field and the dimension may depend on the group.

\textbf{Supplied qualification.} The source asserts that the answer is affirmative for residually finite-soluble groups, citing L. Pyber and D. Segal (2007). Here the conventional meaning of \emph{residually finite-soluble} is that every nonidentity element survives in some finite soluble quotient; a group is soluble if its derived series reaches the trivial subgroup after finitely many steps. The supplied record does not explicitly define this term, and the cited assertion has not been independently verified. The statement begins with an attribution to O. I. Tavgen', whereas the background and proposer metadata name L. Pyber; their respective roles are unresolved in the supplied evidence.
\subsection{Short English statement}
Must every residually finite group that is a product of finitely many cyclic subgroups embed in \(\operatorname{GL}_n(K)\) for some finite \(n\) and some characteristic-zero field \(K\)? The source reports an affirmative answer for residually finite-soluble groups. The registry labels the general question open; this is a supplied status, not an independently verified assessment.
\subsection{Sources}
\begin{itemize}
\item Local UnsolvedMath registry, supplied complete record: permanent ID \(701735\), exact source code \texttt{catalogue entry}, title \emph{Kourovka Notebook Problem 20.77}. The record labels the problem open, lists L. Pyber as proposer, and gives the proposed year as 2022.
\item \emph{The Kourovka Notebook (Unsolved Problems in Group Theory)}, problem 20.77. The supplied background identifies the 20th issue (2022) and the 21st issue edition (2026), with identifier \texttt{arXiv:1401.0300v48}. The collection metadata names E. I. Khukhro and V. D. Mazurov as editors. These bibliographic details are supplied and have not been independently checked.
\item Attribution qualification: the supplied statement opens with ``(O. I. Tavgen')'', while the background and proposer metadata name L. Pyber. The record does not explain the relationship between these attributions.
\item L. Pyber and D. Segal, \emph{J. Reine Angew. Math.} \textbf{612} (2007), 173--211. The source cites this reference for the affirmative residually finite-soluble case; no article title or link is supplied, and the result has not been independently verified.
\end{itemize}
\end{document}
